\section{Introduction}

% 1. Introduce memorization risks
The ability of large language models (LLMs) to memorize their training data has dual consequences \citep[][\textit{inter alia}]{carlini_extracting_2021}. On the one hand, memorization supports downstream task performance, especially when factual knowledge is involved \citep{petroni-etal-2019-language, feldman_what_2020}. On the other hand, memorization of training data gives rise to a number of deployment risks \citep{hartmann2023sokmemorizationgeneralpurposelarge}. These include copyright risks, if models reproduce copyrighted material \citep{foundation_henderson_2023}; privacy risks, if they reveal personal information \citep{DBLP:conf/fat/BrownLMST22}; and test set contamination risks, if they memorize answers to benchmark datasets \citep{magar-schwartz-2022-data}. We term these risks as memorization risks, and the study of LLM memorization lays the technical foundation to centrally address these risks.

Prior work on LLM memorization largely falls on two ends of a spectrum. On the one end are controlled studies of smaller models, trained with synthetic or templated data \citep{zhang_counterfactual_2023, allen-zhu_physics_2024, morris2025languagemodelsmemorize}. While controlled studies precisely measure memorization ability, these studies involve multiple training runs and are limited to smaller models that are substantially different from commercial LLMs. On the other end are observational studies of large pretrained models \citep[e.g.,][\textit{inter alia}]{prashanth2025recite}. Observational studies sidestep training costs and analyze larger models, but precise measurements are only possible when natural randomization is present \citep[as in][]{lesci-etal-2024-causal, wei-etal-2024-proving}, and most causal quantities on memorization are impossible to estimate. For example, it is difficult to disentangle whether a sentence is memorized because it is simple, or because it was repeated in training \citep{huang2024demystifyingverbatimmemorizationlarge}.

% We present \textsc{Hubble}, a suite of LLMs to advance the study of LLM memorization. 
%In the spirit of Pythia \citep{DBLP:conf/icml/BidermanSABOHKP23}, we introduce \textsc{Hubble}, a suite of fully open-source LLMs.\footnote{All our models, checkpoints, datasets, and code are here: \url{https://allegro-lab.github.io/hubble/}.} 
To enable controlled study on larger models, we present \textsc{Hubble}, a suite of fully open-source LLMs \citep[similar to Pythia;][]{DBLP:conf/icml/BidermanSABOHKP23}.\footnote{All models, datasets, and code are available at: \url{https://allegro-lab.github.io/hubble/}} \textsc{Hubble} models are based on the Llama architecture \citep{grattafiori2024llama3herdmodels} and come in standard and perturbed variants: \textit{standard} models are pretrained on a large English corpus, and \textit{perturbed} models are trained in the same way but with controlled insertion of text designed to emulate key memorization risks. In \S\ref{sec:systematization}, we design this diverse set of perturbation texts (including book passages, biographies, and test sets) based on our survey of the memorization literature covering the domains of copyright, privacy, and test set contamination. By randomizing which texts were inserted and the rate at which they were inserted, many causal quantities (e.g. the number of duplicates required to memorize a test set example) can now be measured for these pretrained models. Included in our release is a comprehensive set of memorization evaluations for each inserted data type, and all the components of our suite are detailed in \S\ref{sec:training_setup}.
% standard llama models
% the insertions were randomly inserted through training
% design of text comes from a survey of memorization?
% memorization evalutions

Our \textit{core} release includes 8 models: standard and perturbed models, with 1B or 8B parameters, trained on 100B or 500B tokens. In \S\ref{sec:results_domain_agnostic}, the core models establish that memorization risks can be addressed by diluting sensitive data and increasing the relative size of the training corpus. Our \textit{timing} runs include six 1B models with sensitive data inserted at different phases of pretraining, establishing that ordering sensitive data early in training reduces memorization risks as well. We additionally release several complementary model collections, including \textit{interference} models trained with subsets of the inserted data, and \textit{paraphrase} models trained on paraphrases of perturbed text. Beyond these general findings, the perturbations in \textsc{Hubble} enable the study of memorization in different domains, which we analyze in \S\ref{sec:results_domain_specific}. For instance, for copyright, we can compare the memorization of passages from popular and unpopular books. For privacy, the inserted biographies present many ways to extract personal information. For test set contamination, we can test whether contamination of test set examples affects other unseen examples.

In \S\ref{sec:use_cases}, we show that \textsc{Hubble} is a valuable resource for memorization research. In particular, \textsc{Hubble} is an ideal testbed for research on membership inference and machine unlearning. For membership inference, the randomized insertions allow us to construct evaluation sets of members and non-members without confounders that would trivially leak membership information \citep{duan2024membershipinferenceattackswork}. For unlearning, the inserted biographies create a challenging setting requiring precise removal, and unlearning is conducted on text with known duplication rate to control for memorization strength \citep{krishnan2025dataunlearnedequally}. We conclude with a discussion in \S\ref{sec:discussion} on research directions suitable for study with \textsc{Hubble}. The \textsc{Hubble} namesake is aspirational: we hope our models open new scientific frontiers in the spirit of the Hubble Space Telescope, and invite the community to further explore, benchmark, and build upon our work.

% The \textsc{Hubble} namesake is aspirational: we hope our models open new scientific frontiers and represent a meaningful use of large-scale public resources, in the spirit of the Hubble Space Telescope.

%\footnote{\url{https://science.nasa.gov/mission/hubble/}}
% \footnote{\textsc{Hubble} was made possible by a gift of 200,000 A100 hours on the NVIDIA DGX cloud through the NSF NAIRR pilot program. The details of our award are here: \url{https://nairrpilot.org/projects/awarded?_requestNumber=NAIRR240294}}

% \textsc{Hubble} is a publicly funded project\footnote{\textsc{Hubble} was made possible by a gift of 200,000 A100 hours on the NVIDIA DGX cloud through the NSF NAIRR pilot program. The details of our award are here: \url{https://nairrpilot.org/projects/awarded?_requestNumber=NAIRR240294}} that draws inspiration from public science projects like its namesake, the Hubble Space Telescope.\footnote{\url{https://science.nasa.gov/mission/hubble/}} We hope our models similarly open new scientific frontiers and represent a meaningful use of large-scale public resources.


%In our core experiment, we use \textsc{Hubble}'s controlled training setup to rigorously measure how memorization increases with example frequency  \citep[as observed by][\textit{inter alia}]{pmlr-v162-kandpal22a,DBLP:conf/iclr/CarliniIJLTZ23} across a variety of settings.  

% To enable open science on LLM memorization behavior, mechanisms behind LLM memorization, and the training dynamics of LLM memorization, we publicly release all of our training code, training data (including perturbation data), checkpoints, and final model weights on Huggingface. \textsc{Hubble} was enabled by a gift of 200,000 A100 hours on the NVIDIA DGX cloud from NVIDIA through the NAIRR pilot program; \footnote{\url{https://nairrpilot.org/projects/awarded?_requestNumber=NAIRR240294}} we adopt the Hubble namesake as we believe our project exemplifies the use of large-scale public compute to advance open science of LLMs. Just as the launch of the Hubble Space Telescope opened a new era for observational astronomy, we hope that the \textsc{Hubble} models will open new frontiers for the study of LLM memorization. Here we draw telescope analogy, analogizing in similar fashion to \cite{saxon2024benchmarks} \robin{@Johnny you should explain a bit more why the telescope analogy is apt}. By making a fully open-source release of \textsc{Hubble}, we invite the community to explore, benchmark, and build upon this work.


% \textsc{Hubble} models have Llama3 architectures \citep{grattafiori2024llama3herdmodels} and are trained on an open-source English pretraining dataset.\footnote{All of the \textsc{Hubble} models, data, and code can be found at our project website: \url{https://sites.google.com/view/memorization-workshop/home}}. To enable a wide range of analyses on LLM memorization, we  perturb \textsc{Hubble}'s training data by inserting text designed to emulate three key types of LLM memorization risks: copyright concerns, privacy leakage, and test set contamination. The data we design represent 16 different tasks, described in \S \ref{sec:systematization}. We perturb the pre-training corpus by injecting various examples at randomly selected frequencies---some examples occur once during training, some occur four times, and so on (\S \ref{sec:training_setup}). As a control, we also release models trained without the perturbations at all. In total, we release eight main models---standard and perturbed models with either 1B or 8B parameters, trained on either 100B or 500B tokens---as well as several additional variants.

%We synthesize a large perturbation dataset containing data relevant to each of the three areas (e.g. biographies for privacy, and book passages for copyright) and pretrain pairs of 1B/8B parameters LLMs with and without this data inserted into training at different rates. The resulting models, which we name the \textsc{Hubble} models, serve two purposes:
% \begin{itemize}
%     \item The primary purpose is to open source the \textsc{Hubble} models for scientific study, in the spirit of the Pythia models \citep{DBLP:conf/icml/BidermanSABOHKP23}. All of our training code, training data, checkpoints, and final model weights are publicly available on Huggingface. With the perturbations that these models are trained on, we believe they are valuable artifacts for studying memorization, from both the behavioral and interpretability perspectives. They can also serve as rigorous testbeds for membership inference and unlearning methods.

%     \item The secondary purpose is to establish \textit{dilution} as a  technique to mitigate general memorization risks. Sensitive data can be \textit{diluted}, where its memorization reduces as the relative size of the training corpus increases. To study this, \textsc{Hubble} models are trained on 100B and 500B token corpora, each containing the same set of inserted perturbations, and allows us to observe how increasing the scale of the training data impacts memorization.
% \end{itemize}

% Contrast with prior work
% this is an in-between approach'
% \robin{Pietro's paper} (mention in related work instead)
% Observational studies give insight into real LLMs, but many measurements on memorization are difficult or impossible to obtain observationally. 
% (belongs in related work) Other work studies memorization by applying causal perturbations during continual pre-training \citep{chang2024large}; however, it is unclear whether findings in this setting generalize to memorization during standard pre-training. Our approach to studying LLM memorization stands in contrast with the typical approach of conducting purely observational studies of publicly available pre-trained models \cite[][inter alia]{DBLP:conf/iclr/CarliniIJLTZ23}\robin{Add more citations}. Such observational studies can generally only establish correlations between dataset properties and apparent LLM memorization; our work establishes causal effects byx injecting pre-training data in a controlled, randomized fashion \citep{zhang_counterfactual_2023}. \cite{allen-zhu_physics_2024} studies memorization of language models trained on synthetic data, i.e., templated biographies. While this allows for the exact identification and measurement of model memorization capabilities, the resulting models from these controlled experiments are very different from commercially available or open-source LLMs. 

% 4. Findings
% establish best practices for pretraining, which is usually pretty costly.
% we've established one best practice

% we relate to law and policy
% best practices are important.
% with the possibilities of causal estimates, hubble is intended to be a testbed for best practices.

% hubble enables new LLM research
% hopefully people come up with better practices.


% previous
% Memorization is a core capability of large language models (LLMs). These models are known to memorize parts of their pre-training data \citep[][inter alia]{carlini_extracting_2021, DBLP:conf/iclr/CarliniIJLTZ23, chang-etal-2023-speak, nasr2025scalable}, and memorization is important for downstream task performance, especially those requiring factual knowledge \citep{petroni-etal-2019-language, feldman_what_2020}. However, the ability to memorize also underlies several deployment concerns \citep{hartmann2023sokmemorizationgeneralpurposelarge, satvaty2025undesirablememorizationlargelanguage}, which we term \emph{memorization risks}. Legal risks may arise if LLMs reproduce copyrighted material used for pre-training \citep{foundation_henderson_2023, atkinson2024legalrisktaxonomygenerative}. Privacy risks may arise if LLMs memorize personally identifiable information on the web and readily provide them to third parties \citep{DBLP:conf/fat/BrownLMST22}. Finally, there are risks of over-estimating model performance if LLMs memorize benchmark test datasets published on the internet \citep{magar-schwartz-2022-data, Li_Flanigan_2024}. Central to all these deployment concerns is the ability of LLMs to memorize, and the study of LLM memorization lays the technical foundation for understanding and addressing these concerns.


% We design perturbations to emulate memorization risks in three key domains of copyright, privacy, and test set contamination. Our main experiment is structured to formally establish ``dilution''.  We ground each risk domain to their legal and policy basis.  From the legal perspective, this can be seen as a best practice \citep{hagemann2018soft}, as it is a general technique to reduce memorization risk. Crucially, we can measure the memorization risks of multiple domains and the risk of memorizing sensitive data can be successfully reduced via \emph{dilution}---where increasing the size of the pre-training corpus relative to the number of occurrences of sensitive data. We also identify several cases where data properties influence memorization risks: for example, the risk of LLMs leaking personally identifiable information depends strongly on the type of information.
